\documentclass[../DoAn.tex]{subfiles}
\begin{document}

\section{Kết luận}

Đồ án đã nghiên cứu và triển khai thành công bài toán tìm đường đa tác nhân (MAPF)
trong game bắn xe tăng 2D nhiều người chơi xây dựng trên Unity Engine. Ba cấp độ
thuật toán -- A* baseline, PIBT~C\# và PIBT~C++/TCP -- được triển khai hoàn chỉnh,
chia sẻ cùng tầng điều khiển vật lý và được đánh giá định lượng trên tập bản đồ
chuẩn MovingAI MAPF benchmark.

Về các kết quả chính, đồ án đã xây dựng được hệ thống đọc và dựng bản đồ động từ tệp
\texttt{.map} tại thời điểm chạy, hỗ trợ năm bản đồ kích thước từ 32×32 đến 251×180 mà
không cần chỉnh sửa mã nguồn. Thuật toán A* baseline hoàn thành nhiệm vụ trên bốn trong
năm bản đồ với thời gian ổn định, trong đó Maze128 là giới hạn chung của cả ba thuật toán
do đường đi dài và phức tạp. Hai biến thể PIBT~C\# và PIBT~C++/TCP thể hiện khả năng phối
hợp đa agent không va chạm với thời gian hoàn thành tương đương A* trên bốn trong năm bản
đồ, đồng thời giảm được tình trạng deadlock ở lớp vật lý. Trên nền PIBT~C++/TCP, đồ án còn
cải tiến mô hình xoay của Enemy Agent dựa trên EPIBT, loại bỏ hiện tượng xe tăng xoay giật
$90^\circ$ tại chỗ. Cuối cùng, hệ thống backtest tự động chạy ổn định 900 run (450 tĩnh và
450 động) trên năm mức agent từ 6 đến 72, xuất kết quả CSV và biểu đồ; kết quả xác nhận cả
ba thuật toán duy trì được nhiệm vụ khi môi trường có chướng ngại động (mức tăng replan
30\%--50\%), đồng thời cho thấy thời gian hoàn thành và tỉ lệ timeout tăng dần theo số agent.

So với mục tiêu đề ra tại Mục~\ref{section:1.2}, cả bốn mục tiêu đều đã hoàn thành. So với
các giải pháp tương tự, đóng góp của đồ án nằm ở việc tích hợp và so sánh thực nghiệm ba
cấp độ MAPF trong cùng một môi trường game có vật lý thực tế -- điều chưa thấy trong các
nghiên cứu được khảo sát.

\section{Hạn chế của đồ án}

Bên cạnh các kết quả đạt được, đồ án còn một số hạn chế. Thứ nhất, quy mô thực nghiệm dừng ở 72
agent trên năm bản đồ -- nhỏ hơn nhiều so với quy mô hàng nghìn agent của các benchmark LMAPF
chuyên dụng, nên chưa kiểm chứng được hành vi ở mật độ rất cao. Thứ hai, số lần tái hoạch định
được định nghĩa khác nhau giữa ba thuật toán nên hiện chỉ so sánh được \emph{xu hướng} phía máy
khách, chưa phản ánh chi phí phối hợp bên trong máy chủ C++; bổ sung trường đếm replan thật từ
phía server là cần thiết để so sánh tuyệt đối. Thứ ba, ở các mức agent lớn, một phần số liệu mang
tính ngoại suy theo mô hình thay vì đo trực tiếp đủ số lần lặp, nên cần thêm thực nghiệm để củng
cố độ tin cậy. Thứ tư, cơ chế xét lại agent cùng hai mở rộng LNS và Graph Guidance của EPIBT hiện
được giữ ở mức tối thiểu, nên tiềm năng của chúng trong bối cảnh game chưa được khai thác. Những
hạn chế này đồng thời mở ra các hướng phát triển trình bày dưới đây.

\section{Hướng phát triển}

Trong ngắn hạn, một số việc cần làm để hoàn thiện hệ thống hiện tại. Trước hết là bổ sung
trường \texttt{replan\_count} vào giao thức TCP \texttt{plan\_result} để server C++ trả về
số lần replan thật thay vì phải suy ra từ PIBT~C\#. Tiếp theo là thêm seed tất định cho từng
run để đảm bảo tái lập được hoàn toàn và loại bỏ nhiễu từ quá trình spawn ngẫu nhiên.

Trong dài hạn, đồ án có thể được nâng cấp về cả thuật toán lẫn hệ thống. Hướng đầu tiên
là tích hợp LNS2 \cite{li2021anytime}, một thuật toán MAPF anytime hiệu năng cao cho phép
cải thiện dần chất lượng giải pháp theo thời gian tính toán còn lại. Hướng thứ hai là bổ
sung cơ chế phối hợp tấn công theo đội hình, thay cho việc các agent tấn công Eagle một
cách độc lập như hiện tại, nhằm tăng hiệu quả và tính hấp dẫn của gameplay. Hướng thứ ba
là mở rộng hỗ trợ bản đồ động hoàn toàn, trong đó agent có thể phá hủy tường và thay đổi
đồ thị bản đồ trong thời gian thực, đòi hỏi thuật toán MAPF có khả năng tái hoạch định
toàn cục. Hướng cuối cùng là triển khai WebGL với PIBT server đặt trên cloud thay vì
localhost, mở ra khả năng chạy backtest phân tán và demo trực tuyến không cần cài đặt.

Riêng với nhánh EPIBT, đồ án mới khai thác phần lõi (thao tác đa hành động và kế thừa thao tác)
trong khi giữ cơ chế xét lại agent ở mức tối thiểu và chưa bật hai mở rộng Large Neighborhood
Search và Graph Guidance. Bật và tinh chỉnh các thành phần này -- vốn là yếu tố giúp EPIBT đạt
hiệu năng hàng đầu trên benchmark LMAPF -- là một hướng nâng cấp tự nhiên nhằm cải thiện thêm khả
năng phối hợp và giảm tắc nghẽn ở các bản đồ nhiều hành lang hẹp.

Về khả năng mở rộng quy mô, hệ thống có thể phát triển theo ba trục song song. Trục thứ
nhất là số lượng agent: nâng từ vài agent hiện tại lên hàng chục rồi hàng trăm agent như
quy mô benchmark của LoRR, đòi hỏi tối ưu bộ giải và truyền thông TCP theo lô. Trục thứ
hai là số người chơi đồng thời: mở rộng tầng mạng nhiều người chơi với cơ chế phòng chờ,
ghép trận và đồng bộ trạng thái tin cậy khi nhiều người cùng tham gia một trận. Trục thứ
ba là hạ tầng: triển khai PIBT server trên cloud có cân bằng tải và khả năng co giãn theo
nhu cầu, cho phép phục vụ đồng thời nhiều trận đấu và chạy backtest phân tán khi cả số
agent lẫn số người chơi cùng tăng.

\end{document}
